\documentclass[11pt,a4paper]{article}
\usepackage[margin=2.3cm]{geometry}
\usepackage{amsmath,amssymb,bm}
\usepackage{booktabs,array,longtable}
\usepackage{enumitem}
\usepackage{xcolor}
\usepackage{url}
\usepackage[colorlinks=true,linkcolor=blue!50!black,urlcolor=blue!50!black]{hyperref}
\setlength{\parskip}{0.45em}
\setlength{\parindent}{0pt}
\setlength{\emergencystretch}{3em}
\newcommand{\sig}{\bm{\sigma}}
\newcommand{\Aa}{A_\alpha}
\newcommand{\gfb}{g^{\bar f}}
\newcounter{qn}
\newcommand{\Q}[1]{\refstepcounter{qn}\par\medskip\noindent\textbf{Q\theqn.\ #1}\par\nopagebreak}
\newcommand{\A}{\noindent\textit{Answer.}\ }

\title{Defence Preparation Guide\\[0.3em]
\large What changed since the 21 September draft, why, and how to defend it}
\author{Oyewo Temidayo Solomon\\ Department of Physics, University of Ibadan}
\date{28 September 2026}

\begin{document}
\maketitle
\tableofcontents
\newpage

% =====================================================================
\section*{How to use this guide}
\addcontentsline{toc}{section}{How to use this guide}

Part~\ref{part:changes} lists every substantive change made since the
draft Prof.\ Oyewande reviewed on 21 September, with the old and new
statement side by side. Part~\ref{part:physics} explains the physics and
statistics behind the changes at the depth a panel is likely to probe.
Part~\ref{part:qa} is a bank of likely questions with answers you can
defend. Part~\ref{part:ref} collects the numbers you should know by heart
and tells you where each result lives in the repository.

Read Part~\ref{part:physics} slowly, with the notebooks open. Every
claim in it was checked either against the primary papers or by running
the code; the places where something was \emph{not} re-verified are
flagged explicitly in Section~\ref{sec:unverified}.

% =====================================================================
\section{What changed since 21 September}
\label{part:changes}

\subsection{Theory}

\begin{longtable}{@{}p{0.47\textwidth}p{0.47\textwidth}@{}}
\toprule
\textbf{21 September draft} & \textbf{Current version} \\
\midrule
\endhead
Three coefficients derived: $b_\mu$, $H_{\mu\nu}$, $d_{\mu\nu}$.
& Four: $b_\mu$, $d_{\mu\nu}$, $g_{\lambda\mu\nu}$, $H_{\mu\nu}$, shown to be
the complete set producing spin-dependent terms. \\[0.4em]
$d_{ij}$ and $d_{00}$ ``correct in structure but carry an open overall
sign''; objective counted as four of five achieved.
& Sign resolved: Kostelecký and Lane's lower-index $p_j$ is the covariant
component $p_j=-p^j$. All three $d_{\mu\nu}$ components match exactly.
All five objectives met. \\[0.4em]
Limitation: ``only three of the eight SME coefficients derived''.
& Removed. $a_\mu$, $c_{\mu\nu}$, $e_\mu$ enter only spin-independent
terms; $f_\mu$ does not enter at this order and can be removed by a field
redefinition (Altschul 2006). \\[0.4em]
$e_\mu$, $f_\mu$, $g_{\lambda\mu\nu}$ called ``dimension-5''.
& They are dimensionless. The correct point (Colladay and Kostelecký
1998) is that no renormalisable, gauge-invariant term of the full SME
generates them, so they are expected to be suppressed, though possibly
not for nucleons (Kostelecký and Lane 1999a). \\[0.4em]
Antiparticle Hamiltonian taken as minus the lower block at the same
$\mathbf p$ and $\sig$.
& Antiparticle compared as the CPT-conjugate state, with the momentum
reversed: $H^{\bar f}(\mathbf p,\sig)=-H_{\rm lower}(-\mathbf p,\sig)$.
Verified for every SME coefficient family. \\[0.4em]
Master table last column: ``$|\Aa|\to1$ predicted'' for CPT-odd rows,
``$\Aa\approx0$'' for CPT-even rows.
& Replaced by the actual relation for CPT-conjugate states,
$\gfb=-g^f$ (CPT-odd) or $\gfb=g^f$ (CPT-even). The old entries
contradicted the thesis's own divergence result. \\[0.4em]
Both Kostelecký--Lane papers cited as ``1999''.
& Split into 1999a (PRD 60, 116010) and 1999b (J.\ Math.\ Phys.\ 40,
6245), each cited where used. Altschul (2006) and Wu \textit{et al.}\
(2023) added. \\
\bottomrule
\end{longtable}

\subsection{Database and results}

\begin{longtable}{@{}p{0.47\textwidth}p{0.47\textwidth}@{}}
\toprule
\textbf{21 September draft} & \textbf{Current version} \\
\midrule
\endhead
273 datasets; 21 classified \texttt{UNKNOWN}.
& 283 compiled, 247 analysed. The 36 held out are itemised: 28 combined
envelopes (12 $g_Ag_V$, 16 $g_pg_s$), 5 duplicates at 90\% CL, 2 template
placeholders, 1 bound with no potential assignment. \\[0.4em]
19 of 273 datasets involve antimatter.
& 22 of 247. Muonium ($e^-\mu^+$) bounds moved to the antimatter side,
consistent with how $dd\mu^+$ was already treated. \\[0.4em]
Ten matched pairs.
& Fifteen, of which thirteen are independent (two $g_Vg_V$/$g_pg_p$ rows
use byte-identical Fadeev files). \\[0.4em]
Six of eleven potentials, four of nine sectors with no antimatter data.
& Seven of twelve potentials; six of eleven sector pairs. \\[0.4em]
Priority sectors $e$-$N$, $e$-$n$, $n$-$p$, $n$-$n$.
& $e$-$N$, $n$-$N$, $p$-$N$, $e$-$n$, $n$-$p$, $n$-$n$ (150 of the 225
matter-sector datasets). \\[0.4em]
$g_sg_s$ pair described as a ``CPT-even channel''.
& Described as spin-independent; $a_\mu$ (CPT-odd), $c_{\mu\nu}$ and
$e_\mu$ (CPT-odd) all feed spin-independent terms, so ``CPT-even'' could
not be defended. \\
\bottomrule
\end{longtable}

\subsection{Statistics}

\begin{longtable}{@{}p{0.47\textwidth}p{0.47\textwidth}@{}}
\toprule
\textbf{21 September draft} & \textbf{Current version} \\
\midrule
\endhead
``Every pair remains significant at effectively $p\to0$.''
& Bug fixed: the full $\chi^2$ had been compared with the reduced
degrees of freedom. Now $\chi^2_{\rm eff}=\chi^2 n_{\rm eff}/n$ on
$n_{\rm eff}$ dof, quoted as a Gaussian $Z$. Range $Z=5.1$ to $51$. \\[0.4em]
Bootstrap interval on mean $|\Aa|$ by delta-method, clipped at 1.
& Couplings perturbed log-normally, $\Aa$ recomputed exactly, median
bootstrap bias removed. Every interval now contains its point estimate
(before, six did not). \\
\bottomrule
\end{longtable}

\subsection{Errors found in the final review and corrected}

\begin{itemize}[leftmargin=1.2em]
\item The printed $H_{\rm NR}(g_{[0k]0})$ and $H_{\rm NR}(g_{[kl]j})$ had
the wrong sign: they were Kostelecký and Lane's expressions without the
$p_j=-p^j$ conversion. Correct forms: $+\sig\cdot(\mathbf p\times
\mathbf g_{00})$ and $+\tfrac12\varepsilon_{klm}g_{mlj}p^j\sigma^k$.
\item ``0.9998, the highest asymmetry'': the highest is 1.0000.
\item ``Every Cong 2025 pair sits at the top of the table'': false; replaced
by a like-for-like comparison (all four rose).
\item $e$-$N$ given as 54 in the text but 50 in the table; ``273 rows''
should be 283; ``16--31'' datasets should be 14--31; atlas caption said 25
curves from 10 experiments, the figure shows 27 from 11.
\item Unit-error pairs: three matched pairs touch an affected dataset, not
two.
\item Delaunay \textit{et al.}\ (2017) was listed as an antimatter
experiment; it is hydrogen and helium spectroscopy.
\item The $\bar p$-He $|\alpha|$ bound was attributed to Ficek \textit{et al.}\
(2018); it comes from Salumbides \textit{et al.}\ (2014).
\item The tensor current $\bar\psi\sigma^{\mu\nu}\psi$ was called C-even in
the notes; it is C-odd (checked numerically).
\end{itemize}

% =====================================================================
\section{The physics and statistics, in depth}
\label{part:physics}

\subsection{The minimal SME fermion Lagrangian}

For one Dirac fermion (Kostelecký and Lane 1999a, Eqs.~1--3),
\begin{equation}
\mathcal L=\tfrac12 i\bar\psi\Gamma^\nu\overleftrightarrow{\partial}_\nu\psi
-\bar\psi M\psi,
\end{equation}
\begin{align}
M&=m+a_\mu\gamma^\mu+b_\mu\gamma_5\gamma^\mu+\tfrac12H_{\mu\nu}\sigma^{\mu\nu},\\
\Gamma^\nu&=\gamma^\nu+c^{\mu\nu}\gamma_\mu+d^{\mu\nu}\gamma_5\gamma_\mu
+e^\nu+if^\nu\gamma_5+\tfrac12g^{\lambda\mu\nu}\sigma_{\lambda\mu}.
\end{align}

\textbf{Index counting.} Each coefficient carries the indices of its
Dirac matrix plus, in $\Gamma^\nu$, the derivative index $\nu$. Scalar
and pseudoscalar structures carry none, vector and axial vector one,
the tensor $\sigma$ two. So $e,f$ have one index, $c,d$ two, and $g$
three, which is why it needs a third letter, $\lambda$. It is
antisymmetric in its first two indices because $\sigma^{\lambda\mu}$ is.

\textbf{Dimensions.} ``The parameters appearing in $M$ have dimensions of
mass, while those in $\Gamma$ are dimensionless'' (Kostelecký and Lane
1999a). So $a,b,H$ have mass dimension 1 and $c,d,e,f,g$ are
dimensionless.

\textbf{Why $e,f,g$ are special.} Colladay and Kostelecký (1998, near
their Eq.~29): all renormalisable terms that could generate them ``are
directly incompatible with the electroweak structure''; they could come
only from higher-dimensional operators involving the Higgs field, so they
are expected to be highly suppressed. Kostelecký and Lane (1999a) add
that for protons and neutrons, strong binding ``might generate effective
terms governed by appreciable parameters''. That caveat is the physical
reason studying $g_{\lambda\mu\nu}$ is worthwhile for nucleon sectors.

\textbf{CPT parity from bilinears.} The CPT parities of
$\bar\psi\Gamma\psi$ are $+$ (scalar), $+$ (pseudoscalar), $-$ (vector),
$-$ (axial vector), $+$ (tensor); each derivative adds a $-$. That gives
$a,b$ odd; $H$ even; $c,d$ even; $e,f,g$ odd, exactly Kostelecký and
Lane's classification.

\subsection{The Foldy--Wouthuysen reduction}

Split the Dirac Hamiltonian into even (block-diagonal) and odd
(off-diagonal) parts relative to $\beta$. The FW generator
$S=-i\beta\mathcal O/2m$ removes the odd part order by order:
\begin{equation}
H'\approx\beta m+\mathcal E+\frac{\beta\mathcal O^2}{2m}
-\frac{[\mathcal O,[\mathcal O,\mathcal E]]}{8m^2}+\dots
\end{equation}
An even perturbation survives at order $m^0$ (e.g.\ $b_i$, $H_{ij}$). An
odd perturbation contributes only through its cross term with
$\bm\alpha\cdot\mathbf p$, at order $1/m$ (e.g.\ $b_0$, $H_{0i}$). Kinetic
coefficients multiply $p_\nu$, and when $\nu=0$ that brings in the
energy, $\approx m$, which is why $d_{i0}$ and $g_{[kl]0}$ are
mass-enhanced (order $m^{+1}$).

Coefficients in $\Gamma_0$ first need the field redefinition
$A=1-\tfrac12\gamma^0(\Gamma_0-\gamma_0)$ (Kostelecký and Lane 1999b) so
that the time-derivative term is standard and the Hamiltonian is
Hermitian.

\subsection{The momentum-index convention (the old ``open sign'')}

Kostelecký and Lane's non-relativistic Hamiltonian (1999a, Eq.~4) is
written with lower-index $p_j$. Their intermediate relativistic
Hamiltonian (1999b) has the free term $m\mathcal P_0:=-p_j\gamma^0\gamma^j$.
That equals $-\bm\alpha\cdot\mathbf p$ if $p_j$ is read as the physical
momentum, which would give the wrong free Dirac Hamiltonian. To recover
$+\bm\alpha\cdot\mathbf p+\beta m$, $p_j$ must be the covariant component,
$p_j=-p^j$. Under that reading every $d_{\mu\nu}$ term agrees, signs
included. The mass-enhanced $d_{i0}$ term has no momentum, which is why it
agreed all along while the momentum-linear terms seemed flipped: that
pattern is the fingerprint of an index convention, not of an error. The
$g_{\lambda\mu\nu}$ sector, derived independently, agrees under the same
reading, an independent check.

\subsection{Completeness of the four coefficients}

In Kostelecký and Lane's Hamiltonian to third order (1999b, Eq.~26),
every term multiplied by $\sig$ contains only $b$, $d$, $g$, $H$. The
spin-independent terms contain $a$, $c$, $e$; $f$ appears nowhere. Since
every Dobrescu--Mocioiu potential except $V_1$ is spin-dependent, the four
coefficients are the complete spin-dependent dictionary at these orders.
Altschul (2006) shows $f_\mu$ can be removed entirely by a field
redefinition.

\subsection{Matter versus antimatter: the central subtlety}
\label{sec:antimatter}

After FW reduction the even Hamiltonian has an upper block (particles)
and a lower block (negative-energy solutions).

\textbf{Hole theory.} An antiparticle of momentum $\mathbf p$ and spin
$\mathbf s$ is the \emph{absence} of a negative-energy solution with
momentum $-\mathbf p$ and spin $-\mathbf s$, and its energy is minus that
solution's energy.

\textbf{Two legitimate comparisons.}
\begin{enumerate}[leftmargin=1.5em]
\item \emph{CPT-conjugate state.} CPT maps a particle $(\mathbf p,\mathbf
s)$ to an antiparticle $(\mathbf p,-\mathbf s)$ (C swaps particle and
antiparticle, P reverses $\mathbf p$, T reverses $\mathbf p$ and
$\mathbf s$). Labelling that state by the particle's spin operator:
\begin{equation}
H^{\bar f}(\mathbf p,\sig)=-H_{\rm lower}(-\mathbf p,\sig).
\end{equation}
Under exact CPT symmetry this equals the particle Hamiltonian, so this is
the comparison under which $\Aa=0$. The thesis defines $\gfb$ this way.
\item \emph{Same physical spin.} Writing $H_{\rm lower}=a_0+\mathbf
a\cdot\sig$, charge conjugation ($\chi\to i\sigma_2\chi^*$) gives
\begin{equation}
H^{\bar f}_{\rm same\ spin}(\mathbf p,\sig)=-a_0(-\mathbf p)+\mathbf a(-\mathbf p)\cdot\sig .
\end{equation}
\end{enumerate}

\textbf{The rule.} Between CPT-conjugate states a coefficient flips sign
exactly when it is CPT-odd. At the same physical spin, the relation is
(CPT parity)$\times$(parity of the term under $\mathbf s\to-\mathbf s$), counting even as $+1$ and odd as $-1$, where $+1$ means the same sign and $-1$ a flip. For $b$: $(-1)(-1)=+1$, same sign; for $H$: $(+1)(-1)=-1$, a flip.
The two differ only for spin-linear terms, and that difference is
Kostelecký and Lane's antiparticle rule (1999b, last paragraph):
$a\to-a$, $b\to+b$, $c\to+c$, $d\to-d$, $e\to-e$, $f\to-f$, $g\to+g$,
$H\to-H$.

\textbf{Worked example: $b_i$.} $\gamma^0\gamma_5\gamma^i=-\Sigma^i$ and
$\Sigma^i={\rm diag}(\sigma^i,\sigma^i)$, so both blocks equal
$-\mathbf b\cdot\sig$. With no momentum dependence,
$H^{\bar f}_{\rm CPT}=-(-\mathbf b\cdot\sig)=+\mathbf b\cdot\sig$:
opposite to the particle. At the same physical spin the antiparticle has
$-\mathbf b\cdot\mathbf s$, the same as the particle. Both statements are
true; they compare different states.

\textbf{Why C alone does not do it.} The axial current is C-even (checked:
$C(\gamma_5\gamma^\mu)^TC^{-1}=+\gamma_5\gamma^\mu$), so charge conjugation
of the bilinear gives no sign. The flip is a CPT effect.

\textbf{What went wrong before.} The earlier method took $-H_{\rm
lower}(\mathbf p,\sig)$, reversing neither momentum nor spin. It agrees
with the CPT-conjugate comparison only for terms even in $\mathbf p$. For
$b_i$ that is harmless, which is why the central result survived. For
$g_{[0k]0}$ and $g_{[kl]j}$, which are linear in $\mathbf p$, it produced
the false claim that only one $g$ family flips.

\textbf{External checks.} Bluhm, Kostelecký and Russell (1999): their
hydrogen-to-antihydrogen substitution leaves $b_\mu$ unchanged (the
same-spin rule), and the $b$-dependent 1S--2S shift reverses because ``the
hyperfine states in $\bar{\rm H}$ have opposite positron and antiproton
spin assignments'': the compared states are CPT conjugates.
\texttt{FW\_antiparticle\_all\_coefficients.ipynb} checks both
comparisons for all 16 coefficient families; both agree with CPT parity
and with the published rule in every case.

\subsection{The asymmetry parameter and why it cannot test CPT with bounds}

$\Aa=(g^f-\gfb)/(g^f+\gfb)$, with $\gfb$ for the CPT-conjugate state.

\begin{itemize}[leftmargin=1.2em]
\item Exact CPT symmetry: $\gfb=g^f$, so $\Aa=0$.
\item A CPT-odd coefficient alone: $\gfb=-g^f$, so the denominator
vanishes and $\Aa$ diverges. It does \emph{not} approach 1.
\item Published bounds are positive numbers $|g|\le B$. Put two of them
in: $\Aa=(1-r)/(1+r)$ with $r=B_{\bar f}/B_f$. It depends only on the
ratio of sensitivities. One bound much tighter than the other gives
$|\Aa|\to1$ whatever the physics.
\item Any true pair of couplings below the two bounds is allowed, so the
true $\Aa$ can be anything. The bounds do not constrain it.
\end{itemize}

\textbf{Empirical confirmation.} The four pairs built on the 2025 hydrogen
bounds of Cong \textit{et al.}\ each have an older twin with the same
coupling, potential and antimatter partner. Tightening only the matter
side raised the asymmetry every time: 0.9998$\to$0.9999,
0.8044$\to$1.0000, 0.8236$\to$0.9303, 0.7994$\to$0.9304.

\textbf{Design criterion.} A new antimatter bound is useful for a CPT test
in proportion to how well its sensitivity matches the matter bound, not
how tight it is. That is why the Jiao/Karshenboim pair (0.3336), the least
saturated, is the most informative.

\subsection{Statistics}

\textbf{Grid.} Both curves are interpolated log-linearly onto 300
log-spaced points over their overlap.

\textbf{Uncertainties.} Published exclusion curves carry no error bars.
Each point gets a fractional uncertainty from the local curvature of the
curve in log--log space, between 2\% and 50\%. This is a heuristic, and
the $\chi^2$ values depend on it; the conclusions do not.

\textbf{Weighted $\chi^2$.} $\chi^2=\sum_i (g_{m,i}-g_{a,i})^2/(\sigma_{m,i}^2
+\sigma_{a,i}^2)$ tests $H_0$: the two exclusion curves coincide on the
grid. Rejecting it means the two experiments reached different
sensitivities. It says nothing about the couplings themselves.

\textbf{Autocorrelation.} Adjacent grid points are not independent. The
autocorrelation length $\tau$ is the lag where the normalised
autocorrelation of the $\chi^2$ residuals first drops below $1/e$;
$n_{\rm eff}=n/\tau$. The statistic is rescaled by the same factor,
$\chi^2_{\rm eff}=\chi^2 n_{\rm eff}/n$, and compared with a $\chi^2$
distribution on $n_{\rm eff}$ dof. (The old code compared the full
$\chi^2$ with $n_{\rm eff}$, which overstated significance.) Tail
probabilities are computed in arbitrary precision and quoted as a
one-sided Gaussian $Z$. Example: Jiao/Karshenboim, $\chi^2_{\rm
eff}=48.5$ on 9 dof, $p\approx2\times10^{-7}$, $Z=5.1$.

\textbf{Bootstrap.} 2000 resamples; each perturbs both couplings by
$g\to g\,e^{\sigma\epsilon}$ and recomputes $\Aa$ exactly. $|\Aa|$ is
folded at 0 and saturates at 1, so resampled means sit off the estimate;
the median offset (the bootstrap bias) is subtracted before quoting the
95\% interval. An interval excluding zero is not, on its own, evidence of
asymmetry, because $|\Aa|$ is biased away from zero by construction.

\subsection{The database}

\begin{itemize}[leftmargin=1.2em]
\item Sources: the compilation of Cong \textit{et al.}\ (RMP 2025) and the
primary papers. Filename prefixes give the potential: \texttt{3} means
$V_3$ alone, \texttt{23} means $V_2+V_3$.
\item Units normalised to metres (m, MeV$^{-1}$, eV$^{-1}$, cm, nm, fm).
$|\alpha|$ (gravity-normalised) and $|g|$ bounds are kept apart; typical
values differ by more than thirty orders of magnitude.
\item Muonium: all six $e$-$\mu$ bounds come from muonium (Karshenboim's
$\mu e$ curve is from muonium hyperfine data; Ohayon 2022 and Stadnik 2023
are $e^-\mu^+$). Muonium contains an antimuon, so these are antimatter.
\item Six datasets have implausible ranges up to $10^{52}$~m (two $V_3$,
two $V_{9+10}$, two $V_{15}$). Three matched pairs touch one, but the
tails lie outside every overlap used.
\item 48 datasets are the same bound filed under several couplings; the
coverage figures collapse them to 199 independent bounds.
\end{itemize}

\subsection{What was not re-verified}
\label{sec:unverified}

\begin{itemize}[leftmargin=1.2em]
\item ``Every potential assignment agrees with the source compilation's
records for the 194 datasets they cover.'' This comes from an earlier
check against the Cong \textit{et al.}\ online records (Zenodo
10.5281/zenodo.14572652) and was not repeated.
\item Why the source compilation files the Fadeev \textit{et al.}\ $V_3$
curves under both $g_Vg_V$ and $g_pg_p$. The Fadeev paper treats
pseudovector and pseudoscalar interactions; the $g_Vg_V$ label comes from
the compilation. The thesis only states the fact that the rows are
identical.
\end{itemize}

% =====================================================================
\section{Likely questions and answers}
\label{part:qa}

\subsection{Scope and contribution}

\Q{In one sentence, what does this thesis do?}
\A It builds a complete dictionary from the spin-dependent SME
coefficients to the Dobrescu--Mocioiu potentials, compiles the published
bounds into one standardised database, and shows what a matter--antimatter
comparison of those bounds can and cannot say about CPT symmetry.

\Q{What is new here? Kostelecký and Lane already derived the non-relativistic Hamiltonian.}
\A Three things. The explicit matching of each SME term to the
Dobrescu--Mocioiu basis, together with a proof that $b$, $d$, $g$, $H$ are
the complete spin-dependent set. A standardised, auditable database of 283
bounds with a matching pipeline. And the result that an asymmetry built
from one-sided bounds tracks sensitivity, not CPT status, shown both
analytically and in the data.

\Q{Your main result is a limitation. Why is that a contribution?}
\A Because it is the thing a reader would otherwise get wrong. All fifteen
pairs are highly significant; without the analysis one would be tempted to
read that as a CPT signal. The thesis shows why it is not, turns that into
a concrete design criterion (sensitivity matching), and identifies exactly
which sectors and potentials need antimatter data.

\Q{Why these four coefficients and not the others?}
\A In Kostelecký and Lane's third-order Hamiltonian, only $b$, $d$, $g$, $H$
appear in spin-dependent terms. $a$, $c$, $e$ appear only in
spin-independent terms, and $f$ not at all. All potentials except $V_1$
are spin-dependent, so nothing is missing.

\subsection{The SME--DM mapping}

\Q{SME coefficients are fixed background fields; DM potentials come from
boson exchange between two fermions. How can one map to the other?}
\A The dictionary matches operator structure, not mechanism. A background
$b_i$ gives a single-particle term $-\mathbf b\cdot\sig$; a boson exchange
gives two-body terms with the same spin structure at each vertex. The
mapping tells you which Dobrescu--Mocioiu structure an experiment must be
sensitive to in order to constrain a given SME coefficient combination. It
does not claim the two are the same physics. This is a fair question and
the honest answer is exactly that.

\Q{Why the Foldy--Wouthuysen transformation?}
\A It gives the non-relativistic Hamiltonian as an explicit $2\times2$
Pauli operator, order by order in $1/m$, which is exactly the form needed
to compare term by term with the potentials. It is also the method
Kostelecký and Lane used, which made a term-by-term check possible.

\Q{What was the $d_{ij}$ sign problem, and how do you know it is solved and not hidden?}
\A The momentum-linear terms appeared to have the opposite sign to
Kostelecký and Lane. Their own free-particle term,
$m\mathcal P_0=-p_j\gamma^0\gamma^j$, only reproduces the correct free
Dirac Hamiltonian if $p_j=-p^j$. With that, all terms agree. Two checks
make it convincing: the mass term, which has no momentum, always agreed;
and the $g$ sector, derived separately, agrees under the same convention.
A wavefunction-renormalisation explanation was tested and ruled out.

\Q{Why does $g$ have three indices? Why $\lambda$?}
\A It multiplies $\sigma_{\lambda\mu}$ (two indices) inside $\Gamma^\nu$
(one more, the derivative index). $\lambda$ is just the next free letter;
it is the SME convention. Note that the thesis also uses $\lambda$ for the
interaction range; context distinguishes them.

\Q{Is $g_{\lambda\mu\nu}$ really dimension five?}
\A No; that was an error in the earlier draft. It is dimensionless and
multiplies a dimension-four operator. What is true is that it cannot come
from a renormalisable, gauge-invariant term of the full SME, so it is
expected to be suppressed for electrons (Colladay and Kostelecký 1998),
but possibly not for nucleons (Kostelecký and Lane 1999a).

\Q{Which $g$ families contribute, and why does $g_{[0k]j}$ vanish?}
\A $g_{[kl]0}$ (Zeeman-like, mass-enhanced, $V_2$), $g_{[0k]0}$
(spin--velocity, $V_7$), $g_{[kl]j}$ (velocity-dependent, $V_7,V_8$).
$g_{[0k]j}$ gives nothing at the orders kept; in Eq.~26 of Kostelecký and
Lane it first appears at order $p^2/m^2$.

\subsection{CPT and antimatter}

\Q{Does $b_\mu$ flip sign between matter and antimatter or not? Kostelecký
and Lane say $b\to+b$.}
\A Both are right; they compare different states. Between CPT-conjugate
states (antiparticle spin reversed) the $b$ term flips, which is the
CPT-odd signature. At the same physical spin it does not, which is the
Kostelecký--Lane rule. They differ because $\mathbf b\cdot\sig$ changes
sign when the spin is reversed. The thesis uses the CPT-conjugate
comparison because that is the one under which exact CPT symmetry gives
$\Aa=0$.

\Q{Why not just use charge conjugation?}
\A The axial current is C-even, so C alone gives no sign difference for
$b$. The difference is a CPT effect. Charge conjugation of the two-spinor
is used, but together with the hole-theory reversal of momentum and spin.

\Q{Your earlier draft said only one $g$ family flips. What happened?}
\A The antiparticle was taken as minus the lower block without reversing
the momentum. That is harmless for terms with no momentum, such as $b_i$,
but wrong for the two momentum-linear $g$ families. With the reversal all
three flip, as a CPT-odd coefficient must. I then checked every SME
coefficient family both ways in one notebook.

\Q{How do you know your antiparticle rule is right?}
\A Three independent checks agree: the full symbolic FW reduction with
charge conjugation for all 16 families; Kostelecký and Lane's published
antiparticle substitution rule; and Bluhm, Kostelecký and Russell's
antihydrogen analysis. A direct numerical diagonalisation of the Dirac
Hamiltonian also gives the same signs.

\Q{Muonium: why is it antimatter?}
\A Muonium is $e^-\mu^+$; the $\mu^+$ is an antimuon. The database already
counted $dd\mu^+$ as antimatter for exactly that reason, so treating
muonium as matter was inconsistent.

\subsection{The asymmetry parameter}

\Q{If $\Aa$ built from bounds cannot test CPT, why define it at all?}
\A It is the natural observable, and with signed measurements of the same
coupling in both sectors it would work. With bounds, it becomes a
diagnostic of sensitivity matching. Showing precisely where the obvious
observable breaks down is part of the result.

\Q{What does a CPT-odd coefficient predict for $\Aa$?}
\A For signed couplings, $\gfb=-g^f$ and $\Aa$ diverges. For magnitudes it
gives $|\gfb|=|g^f|$. Neither gives $|\Aa|\to1$. Values near 1 in the data
come from sensitivity gaps.

\Q{Then why is the lowest-asymmetry pair the most informative?}
\A Because its two bounds are closest in sensitivity, so it is the only
comparison not dominated by the sensitivity ratio.

\Q{The $g_sg_s$ pair: is it CPT-even?}
\A I do not claim that. $V_1$ is spin-independent; $a_\mu$ and $e_\mu$
(CPT-odd) and $c_{\mu\nu}$ (CPT-even) all enter spin-independent terms.
What the pair shows is that a channel none of the four spin-dependent
coefficients feeds has an asymmetry comparable to the spin-dependent
pairs, again pointing to sensitivity, not physics.

\subsection{Statistics}

\Q{What is your null hypothesis?}
\A That the matter and antimatter exclusion curves coincide at every grid
point, tested jointly with a weighted $\chi^2$. Rejecting it means the
experiments reached different sensitivities. It is not a test of the
underlying couplings.

\Q{What does a $p$-value mean for an asymmetry built from upper limits?}
\A Only this: if the two curves were the same curve plus noise of the
assumed size, how unlikely is the observed disagreement. It measures a
sensitivity difference.

\Q{Where do your uncertainties come from? Published bounds have none.}
\A From a curvature heuristic, 2\% to 50\% of the bound. That is a
modelling choice. The $Z$ values depend on it; the qualitative conclusions
do not, because they rest on the analytic argument and on the
like-for-like comparisons, not on the $\chi^2$.

\Q{Why correct for autocorrelation, and how?}
\A 300 interpolated points come from far fewer real measurements, so they
are correlated. I estimate $\tau$ from the residual autocorrelation, set
$n_{\rm eff}=n/\tau$, and rescale $\chi^2$ by the same factor. An earlier
version rescaled only the degrees of freedom, which overstated
significance; that is fixed.

\Q{Your earlier draft said $p\to0$. Now?}
\A $Z$ from 5.1 (Jiao/Karshenboim, $p\approx2\times10^{-7}$) to 51. The
earlier ``$p\to0$'' was partly the bug and partly floating-point underflow.

\Q{Your bootstrap interval excludes zero. Is that evidence of asymmetry?}
\A No. $|\Aa|$ is biased away from zero by the absolute value, so its
interval excludes zero almost always. The interval only describes spread
from the assumed uncertainties. I also correct the median bootstrap bias,
which otherwise pushes intervals off the estimate near 0 and 1.

\Q{Fifteen pairs, but are they independent?}
\A Thirteen are. Two $g_Vg_V$/$g_pg_p$ rows use byte-identical files, and
several pairs share one side (Ficek 2018 is the antimatter side of seven).
The shared datasets mean the pairs are not fully independent even among
the thirteen, which is another reason not to combine them into a single
significance.

\subsection{Data}

\Q{How were potentials assigned?}
\A From filename prefixes (the numeric prefix is the potential number),
explicit tokens, and a small override table, checked against the source
compilation's records. \emph{Caveat:} the claim of agreement for the 194
datasets those records cover comes from an earlier check I did not repeat
before this defence.

\Q{Why exclude 36 datasets?}
\A 28 are combined envelopes over several experiments with no single
potential; 5 are 90\% CL duplicates of 95\% CL curves; 2 are template
placeholders; 1 has no potential that can be established. Each is recorded
with its reason; none is silently dropped.

\Q{What about the implausible interaction ranges?}
\A Six datasets reach up to $10^{52}$~m, almost certainly a mis-detected
unit. Three matched pairs touch one, but the asymmetry is computed only on
the overlap, which in all three cases lies far inside the plausible range.

\Q{Why is the $\bar p$-$N$ bound not used?}
\A It is published as a gravity-normalised strength $|\alpha|$, not a
coupling $|g|$; comparing it with the $p$-$N$ data would need a conversion
the framework does not perform.

\subsection{Limitations and future work}

\Q{What are the main limitations?}
\A The asymmetry cannot distinguish CPT violation from sensitivity gaps
with bounds; the uncertainty model is heuristic; derivations stop at the
orders kept in $1/m$; 36 datasets are held out; antimatter coverage is
thin, which is the finding itself.

\Q{What would turn this into a real CPT test?}
\A Antimatter-sector bounds of sensitivity comparable to their matter
counterparts, or signed measurements. The highest-value sectors are
$e$-$N$, $n$-$N$ and $p$-$N$, and the potentials $V_{1a}$, $V_{9+10}$,
$V_{12+13}$. Reanalysing existing BASE, ALPHA and co-magnetometer data, as
Wu \textit{et al.}\ (2023) did for Lorentz-violation data, is the cheapest
first step.

\Q{What would you do next with more time?}
\A Carry the reductions to higher order in $1/m$; decompose the 28
envelope curves; re-verify all potential assignments against the source
records; add a conversion for $|\alpha|$ bounds; and propagate the curve
uncertainties more rigorously than the curvature heuristic.

\subsection{Awkward questions}

\Q{Did you make mistakes along the way?}
\A Yes, and they are documented: the open $d$ sign, the antiparticle
method, a wrong dimension label, a statistics bug and a mis-classified
sector. Each was found by checking against primary sources or by rerunning
the code, and each fix is recorded with the check that confirms it.

\Q{Is your software trustworthy?}
\A It has an automated test suite (140 tests), including regression tests
for each bug fixed, and the full analysis reruns from one command. The
derivations are in executable SymPy notebooks that assert their own
results against Kostelecký and Lane.

\Q{Could a referee reproduce Table II from your repository?}
\A Yes: run the pipeline, and \texttt{results/tables/asymmetry\_summary.csv}
contains every number in the table, including $n_{\rm eff}$, $Z$ and the
intervals.

% =====================================================================
\section{Quick reference}
\label{part:ref}

\subsection{Numbers to know}

\begin{center}
\begin{tabular}{@{}lr@{}}
\toprule
Datasets compiled / analysed / held out & 283 / 247 / 36 \\
Antimatter / matter datasets (analysed) & 22 / 225 \\
Antimatter by sector & $e\bar p$ 9, $e\bar\mu$ 6, $ee^+$ 5, $\bar p$He 1, $dd\mu^+$ 1 \\
Matched pairs (independent) & 15 (13) \\
Pairs in $V_2$ or $V_3$ & 14 of 15 \\
Mean $|\Aa|$ range & 0.3336 to 1.0000 \\
$n_{\rm eff}$ range & 6 to 21 (of 300) \\
$Z$ range & 5.1 to 51 \\
Potentials with no antimatter data & 7 of 12 \\
Sector pairs with no antimatter data & 6 of 11 \\
$e$-$N$ / $n$-$N$ / $p$-$N$ matter datasets & 50 / 49 / 18 (117 of 225) \\
Independent bounds after de-duplication & 199 \\
Implausible-range datasets & 6 \\
\bottomrule
\end{tabular}
\end{center}

\subsection{The fifteen pairs}

\begin{center}\small
\begin{tabular}{@{}lllllrrr@{}}
\toprule
Coupling & $V$ & Sector & Matter & Antimatter & $|\Aa|$ & $n_{\rm eff}$ & $Z$ \\
\midrule
$g_Ag_A$ & 3 & $e\bar p$ & Cong 2025 & Fadeev 2022 & 1.0000 & 8 & 51.1 \\
$g_Ag_A$ & 2 & $e\bar p$ & Cong 2025 & Ficek 2018 & 0.9999 & 21 & 34.4 \\
$g_Ag_A$ & 2 & $e\bar p$ & Karshenboim 2011 & Ficek 2018 & 0.9998 & 17 & 38.8 \\
$g_Ag_A$ & 2 & $ee^+$ & Ficek 2017 & Karshenboim 2011 & 0.9892 & 21 & 48.2 \\
$g_Ag_A$ & 3 & $ee^+$ & Ficek 2017 & Fadeev 2022 & 0.9539 & 6 & 38.5 \\
$g_Vg_V$ & 3 & $ee^+$ & Fadeev 2022 & Fadeev 2022 & 0.9535 & 6 & 37.1 \\
$g_pg_p$ & 3 & $ee^+$ & Fadeev 2022 & Fadeev 2022 & 0.9535 & 6 & 37.1 \\
$g_pg_p$ & 3 & $e\bar p$ & Cong 2025 & Ficek 2018 & 0.9304 & 6 & 34.3 \\
$g_Ag_A$ & 3 & $e\bar p$ & Cong 2025 & Ficek 2018 & 0.9303 & 6 & 34.3 \\
$g_sg_s$ & 1 & $ee^+$ & Delaunay 2017 & Adkins 2022 & 0.8727 & 21 & 25.7 \\
$g_Ag_A$ & 3 & $e\bar p$ & Fadeev 2022 & Ficek 2018 & 0.8236 & 6 & 37.3 \\
$g_Ag_A$ & 3 & $e\bar p$ & Fadeev 2022 & Fadeev 2022 & 0.8044 & 8 & 43.8 \\
$g_Vg_V$ & 3 & $e\bar p$ & Fadeev 2022 & Ficek 2018 & 0.7994 & 7 & 31.8 \\
$g_pg_p$ & 3 & $e\bar p$ & Fadeev 2022 & Ficek 2018 & 0.7994 & 7 & 31.8 \\
$g_Ag_A$ & 2 & $ee^+$ & Jiao 2020 & Karshenboim 2011 & 0.3336 & 9 & 5.1 \\
\bottomrule
\end{tabular}
\end{center}

\subsection{Matter--antimatter relations}

\begin{center}
\begin{tabular}{@{}lccc@{}}
\toprule
Coefficient & CPT & CPT-conjugate state & Same physical spin \\
\midrule
$b_\mu$ & odd & flips & same \\
$g_{\lambda\mu\nu}$ (all 3 families) & odd & flips & same \\
$H_{\mu\nu}$ & even & same & flips \\
$d_{\mu\nu}$ & even & same & flips \\
$a_\mu$, $e_\mu$ (spin-independent) & odd & flips & flips \\
$c_{\mu\nu}$ (spin-independent) & even & same & same \\
\bottomrule
\end{tabular}
\end{center}

\subsection{Where things live}

\begin{itemize}[leftmargin=1.2em]
\item Derivations: \path{derivations/sympy/FW_bmu_term.ipynb},
\path{FW_Hmunu_term.ipynb}, \path{FW_dmunu_term.ipynb},
\path{FW_gmunu_term.ipynb}.
\item Antiparticle check for every coefficient:
\path{derivations/sympy/FW_antiparticle_all_coefficients.ipynb}.
\item Theory notes: \path{docs/theory_notes/}.
\item Pipeline: \path{spindep_framework}; results in
\path{results/tables/asymmetry_summary.csv} and
\path{dataset_registry.csv}; statistics in
\path{spindep/src/statistics.py}; tests in \path{tests/}.
\item Thesis: \path{thesis/project/}; journal draft:
\path{thesis/aps-draft/}; summit abstract:
\path{thesis/aps-draft/summit-abstract.txt}.
\end{itemize}

\subsection{Key references}

\begin{itemize}[leftmargin=1.2em]
\item D.~Colladay and V.~A. Kostelecký, Phys.\ Rev.\ D \textbf{58}, 116002 (1998).
\item V.~A. Kostelecký and C.~D. Lane, Phys.\ Rev.\ D \textbf{60}, 116010 (1999) [1999a].
\item V.~A. Kostelecký and C.~D. Lane, J.\ Math.\ Phys.\ \textbf{40}, 6245 (1999) [1999b].
\item R.~Bluhm, V.~A. Kostelecký and N.~Russell, Phys.\ Rev.\ Lett.\ \textbf{82}, 2254 (1999).
\item B.~Altschul, J.\ Phys.\ A \textbf{39}, 13757 (2006).
\item B.~A. Dobrescu and I.~Mocioiu, JHEP \textbf{11}, 005 (2006).
\item L.~Cong \textit{et al.}, Rev.\ Mod.\ Phys.\ \textbf{97}, 025005 (2025).
\item S.~G. Karshenboim, Phys.\ Rev.\ A \textbf{83}, 062119 (2011).
\item L.~Y. Wu \textit{et al.}, Phys.\ Rev.\ Lett.\ \textbf{131}, 091002 (2023).
\end{itemize}

\end{document}
